\section{Related work}
\label{sec:related-work}

This section positions the present heuristic within three adjacent literatures: graph algorithms for cycle breaking, experimental algorithmics and heuristic design for combinatorial optimization, and the specific minimal/stable and SCC-aware feedback-set lineage that motivates the proposed component-local refinement step.

\subsection{Feedback arc set algorithms}

Minimum feedback arc set sits between graph-theoretic cycle breaking and ordering optimization. Classical formulations treat it as deleting a minimum-weight arc set to obtain an acyclic digraph or, equivalently, as finding an ordering of minimum backward weight \cite{LC66,F90}. The decision version is NP-hard \cite{K72}, and exact methods remain important mainly as certification tools on restricted instance classes. Modern exact work includes sparse-MFAS formulations and computational exact algorithms \cite{BSNA21} as well as subsequent commentary on formulation and separation choices \cite{GJR22}. Those exact methods motivate the validation subsets in the present paper, but they do not remove the practical need for heuristics on larger sparse graphs -- the territory of experimental algorithmics and algorithm engineering that this paper occupies.

Approximation and reduction-based work provides methodological context rather than the main claim of this paper. The general local-ratio framework of \cite{BYGR98} explains why iterative weight reductions can preserve improvement structure in weighted optimization. For directed feedback problems, Demetrescu and Finocchi \cite{DF03} show how cycle-based local-ratio reductions can be used algorithmically; that work is the main conceptual ancestor of the LR-TA seed used here. Earlier directed-feedback approximation results \cite{ENSS98} and more recent localization-oriented engineering work such as Tight Localizations of Feedback Sets (TIGHT-CUT/TIGHT-CUT*) \cite{HGH21} place the present heuristic study within a broader sparse feedback-set literature.

Classical heuristic baselines remain relevant because many practical implementations still inherit their design logic. Greedy score-based ordering rules in the Eades family \cite{ELS93,EL95} are fast and interpretable, but they act mainly through global score differences rather than targeted treatment of residual cyclic subgraphs. Sorting-heuristic studies on sparse graphs \cite{BH11TR,BH13,Hanauer2017thesis} show that SCC-aware decomposition and sparse-ordering structure can matter materially once the graph is incomplete rather than tournament-like. At larger scales, deployment-oriented engineering becomes central \cite{SST16}, and PageRank-based FAS heuristics provide a more recent general-digraph comparator family \cite{GLT23}.

\subsection{Ordering and ranking formulations}

The ordering interpretation of MWFAS connects the problem to ranking and linear ordering. In dense pairwise-comparison settings, the natural input is often a nonnegative matrix \(W\) whose entry \(W_{ij}\) expresses evidence that item \(i\) should precede item \(j\), and the objective is to maximize total forward weight or minimize total backward weight \cite{ACN08,CFR10,VahidiKoutis2024arxiv}. That viewpoint is especially natural for complete tournaments and dense linear-ordering benchmarks. In sparse digraphs, many ordered pairs are absent altogether, so the problem is more naturally expressed as cycle breaking in an incomplete directed graph. Hanauer's sparse-ordering thesis makes that distinction explicit and studies structural and algorithmic consequences for sparse instances \cite{Hanauer2017thesis}.

This distinction matters for baseline selection and claim scope. The python-igraph interface \cite{python_igraph_feedback_arc_set} provides an accessible graph-library comparator. The DRMacIver/FAS tool \cite{drmaciver_feedback_arc_set}, by contrast, is an external matrix-based ordering heuristic whose documented input model is pairwise-ordering data rather than sparse graph analysis. The evaluated executable exhibits run-to-run variation under internal initialization, so it is treated empirically as stochastic here. This makes it both a useful cross-paradigm comparator on sparse instances and an especially appropriate comparator on dense LOLIB instances, where the matrix-based formulation aligns naturally with the benchmark structure.

The dense linear-ordering literature provides the right background for the LOLIB transfer test. LOLIB was established as a benchmark family for complete weighted ordering problems rather than sparse general digraphs \cite{MRD12,lolib_library}. Tournament and dense-ordering heuristics therefore form a separate comparison class, including weighted local-search methods \cite{FLRS10} and more recent dense-ordering work such as \cite{OZ24}. The present paper uses LOLIB to mark a scope boundary rather than to claim superiority in that dense regime.

\subsection{Minimal/stable FAS and SCC-based heuristics}

SCC structure is a recurring theme in practical feedback arc heuristics because cycles are localized there. Decomposition along SCC boundaries is therefore a natural way to avoid spending effort on regions that are already acyclic. The present work uses that insight twice: in the attributed seed constructions and again in the IPSNS refinement layer, where only SCCs with positive incumbent backward contribution are eligible for repair. The minimal-feedback literature is also relevant here: Cavallaro, Cutello, and Pavone study heuristics for finding small minimal feedback arc sets even on large graphs \cite{CCP23CEUR}, and later minimal/stable work develops weighted variants \cite{CCP24,CC25}. These methods are close in spirit because they combine SCC-aware processing with acyclicity-preserving restoration, but their primary target is minimal or minimal/stable feedback sets rather than the present sparse weighted-MWFAS refinement layer.

The second seed in the framework comes from the minimal/stable feedback arc set line of Cavallaro and collaborators. Cavallaro and Cutello \cite{CC25} study weighted minimal-and-stable feedback arc sets, extending earlier minimal/stable constructions \cite{CCP24}. That lineage matters here because it provides a complementary constructive bias to local-ratio seeding and because it distinguishes \emph{minimal} from \emph{minimum-weight}: an inclusion-minimal feedback arc set is one from which no arc can be restored while preserving acyclicity, but it need not minimize total weight.

This distinction is easy to blur in heuristic writing, so the paper keeps it explicit. LR-TA and the WMSF-style seed both use acyclicity-preserving add-back, and IPSNS inherits the same add-back logic inside SCC-restricted repair. What matters scientifically is not a claim of exact solution, but that these routines provide feasible and complementary seeds for later incumbent-protected refinement.

\subsection{Relationship to preliminary author work and present contribution}
\label{subsec:prior-work}

The public preprint of Vahidi and Koutis \cite{VahidiKoutis2024arxiv} introduced the ranking-oriented framing of MWFAS together with an earlier local-ratio-based pipeline built around cycle peeling, feasible reinsertion, and topological extraction. The present manuscript does not claim that foundation as new. The LR-TA seed descends from that earlier pipeline and is attributed accordingly. The WMSF-style seed is likewise not presented as a new contribution of this paper; it is an adapted implementation of the weighted minimal/stable line of Cavallaro and Cutello \cite{CC25}.

What is new here is the IPSNS-centered framework and the accompanying computational study: dual-seed initialization, supporting feasibility and monotonicity analysis, exact and MIP validation, sparse benchmark comparisons, repeated-run analysis against the evaluated DRMacIver/FAS executable, holdout-supported defaults, and a dense LOLIB scope study. The contribution is therefore a sparse-digraph algorithm-engineering paper centered on incumbent-protected SCC-local refinement, with inherited components attributed rather than rebranded and with dense ranking work treated as related but distinct context rather than as the main target regime.
